\section*{Capítulo \ref{ch:g1:hot-and-cold} --- Quente e frio}

\begin{solution}{exo:g1:hot-and-cold:1}
O cubo de gelo, depois o ar do quarto, depois a água do banho, depois
a sopa recém-tirada da panela.
\end{solution}

\begin{solution}{exo:g1:hot-and-cold:2}
A água da torneira é \emph{fria} (ou mais fria) em relação ao
chocolate quente, mas \emph{quente} (ou mais quente) em relação ao
suco.
\end{solution}

\begin{solution}{exo:g1:hot-and-cold:3}
Por exemplo: o Sol, uma chama, um aquecedor, o seu próprio corpo (ou
as mãos esfregando). O Sol esquenta a praça inteira.
\end{solution}

\begin{solution}{exo:g1:hot-and-cold:4}
As palmas parecem quentes contra o rosto: esfregar esquentou as mãos.
\end{solution}

\begin{solution}{exo:g1:hot-and-cold:5}
O cubo que está no sol derrete primeiro: o Sol o esquenta, e é o
calor que derrete o gelo. O que está na sombra recebe menos calor e
dura mais.
\end{solution}

\begin{solution}{exo:g1:hot-and-cold:6}
A poça é água --- a mesma água que era o cubo de gelo. O gelo duro
não desapareceu; o calor da sala o derreteu e o transformou em água
líquida.
\end{solution}

\begin{solution}{exo:g1:hot-and-cold:7}
Depois de uma manhã inteira na mesma sombra, o escorregador e o banco
estão igualmente frios. O metal só \emph{parece} mais frio para a mão
--- o tato não é um juiz justo, como mostraram a colher de metal e a
colher de pau.
\end{solution}

\begin{solution}{exo:g1:hot-and-cold:8}
Toque numa tigela e logo em seguida na outra, com a \emph{mesma} mão,
e peça a um colega que confira sem escutar a sua resposta. A mesma
mão é necessária porque o truque das três tigelas mostrou que duas
mãos podem discordar sobre a mesmíssima água.
\end{solution}

\begin{solution}{exo:g1:hot-and-cold:9}
Lea está certa. Um casaco não fabrica calor: ele só impede que o
calor atravesse. O boneco de neve não tem calor nenhum para guardar
--- em compensação, o casaco mantém do lado de fora o \emph{ar quente
do dia}, de modo que o boneco enrolado derrete mais devagar e dura
mais tempo.
\end{solution}
